\documentclass[10pt,a4paper]{amsart}
\usepackage{amsmath,amssymb,amsthm,graphicx}
% Retained minimal preamble for the Stacks Project chapter sites-cohomology.tex (Cohomology on Sites).
% Verbatim excerpt of the chapter preamble: Xy-pic setup (native preamble lines 14-18) and the native
% theorem declarations with the mathematical macros (native preamble lines 154-217). Copied mechanically;
% nothing added, removed or redefined. The wrapper class and the amsmath/amssymb/amsthm/graphicx base
% packages are supplied by the assembler, so the class/externaldocument part of the native preamble is omitted.
% For commutative diagrams we use Xy-pic
\usepackage[all]{xy}

% We use 2cell for 2-commutative diagrams.
\xyoption{2cell}
% Theorem environments.
%
\theoremstyle{plain}
\newtheorem{theorem}[subsection]{Theorem}
\newtheorem{proposition}[subsection]{Proposition}
\newtheorem{lemma}[subsection]{Lemma}

\theoremstyle{definition}
\newtheorem{definition}[subsection]{Definition}
\newtheorem{example}[subsection]{Example}
\newtheorem{exercise}[subsection]{Exercise}
\newtheorem{situation}[subsection]{Situation}

\theoremstyle{remark}
\newtheorem{remark}[subsection]{Remark}
\newtheorem{remarks}[subsection]{Remarks}

\numberwithin{equation}{subsection}

% Macros
%
\def\lim{\mathop{\mathrm{lim}}\nolimits}
\def\colim{\mathop{\mathrm{colim}}\nolimits}
\def\Spec{\mathop{\mathrm{Spec}}}
\def\Hom{\mathop{\mathrm{Hom}}\nolimits}
\def\Ext{\mathop{\mathrm{Ext}}\nolimits}
\def\SheafHom{\mathop{\mathcal{H}\!\mathit{om}}\nolimits}
\def\SheafExt{\mathop{\mathcal{E}\!\mathit{xt}}\nolimits}
\def\Sch{\mathit{Sch}}
\def\Mor{\mathop{\mathrm{Mor}}\nolimits}
\def\Ob{\mathop{\mathrm{Ob}}\nolimits}
\def\Sh{\mathop{\mathit{Sh}}\nolimits}
\def\NL{\mathop{N\!L}\nolimits}
\def\CH{\mathop{\mathrm{CH}}\nolimits}
\def\proetale{{pro\text{-}\acute{e}tale}}
\def\etale{{\acute{e}tale}}
\def\QCoh{\mathit{QCoh}}
\def\Ker{\mathop{\mathrm{Ker}}}
\def\Im{\mathop{\mathrm{Im}}}
\def\Coker{\mathop{\mathrm{Coker}}}
\def\Coim{\mathop{\mathrm{Coim}}}

% Boxtimes
%
\DeclareMathSymbol{\boxtimes}{\mathbin}{AMSa}{"02}

%
% Macros for moduli stacks/spaces
%
\def\QCohstack{\mathcal{QC}\!\mathit{oh}}
\def\Cohstack{\mathcal{C}\!\mathit{oh}}
\def\Spacesstack{\mathcal{S}\!\mathit{paces}}
\def\Quotfunctor{\mathrm{Quot}}
\def\Hilbfunctor{\mathrm{Hilb}}
\def\Curvesstack{\mathcal{C}\!\mathit{urves}}
\def\Polarizedstack{\mathcal{P}\!\mathit{olarized}}
\def\Complexesstack{\mathcal{C}\!\mathit{omplexes}}
% \Pic is the operator that assigns to X its picard group, usage \Pic(X)
% \Picardstack_{X/B} denotes the Picard stack of X over B
% \Picardfunctor_{X/B} denotes the Picard functor of X over B
\def\Pic{\mathop{\mathrm{Pic}}\nolimits}
\def\Picardstack{\mathcal{P}\!\mathit{ic}}
\def\Picardfunctor{\mathrm{Pic}}
\def\Deformationcategory{\mathcal{D}\!\mathit{ef}}

% --- appended 2026-10-12 (Chief): the project own definitions of the reference/slogan/history
% display environments, copied mechanically from the shared Stacks preamble lines 8-12
% (\usepackage{verbatim} plus the three \newenvironment lines, which the chapter preamble keeps
% verbatim from the project preamble). Nothing added, removed or redefined: these environments
% carry no mathematical content and the project itself typesets them as comments. Without these
% three lines a counted statement whose own text uses them fails with Environment slogan undefined.
\usepackage{verbatim}
\newenvironment{reference}{\comment}{\endcomment}
\newenvironment{slogan}{\comment}{\endcomment}
\newenvironment{history}{\comment}{\endcomment}
\usepackage{cleveref}
\crefname{originalequation}{original equation}{original equations}
\Crefname{originalequation}{Original equation}{Original equations}
\makeatletter
\newlabel{equation-cohomology-hom-complex}{{0A5X}{source}{}{sourcepointer-1}{}}
\newlabel{equation-cohomology-hom-complex@cref}{{[originalequation][0][]0A5X}{[1][0][]source}{}{}{}}
\makeatother
\crefname{originalequation}{original equation}{original equations}
\Crefname{originalequation}{Original equation}{Original equations}
\makeatletter
\newlabel{equation-global-cohomology-hom-complex}{{0A8Z}{source}{}{sourcepointer-2}{}}
\newlabel{equation-global-cohomology-hom-complex@cref}{{[originalequation][0][]0A8Z}{[1][0][]source}{}{}{}}
\makeatother
\crefname{originalreference}{original reference}{original references}
\Crefname{originalreference}{Original reference}{Original references}
\makeatletter
\newlabel{lemma-restrict-K-injective-to-open}{{08BS}{source}{}{sourcepointer-3}{}}
\newlabel{lemma-restrict-K-injective-to-open@cref}{{[originalreference][0][]08BS}{[1][0][]source}{}{}{}}
\makeatother
\title{The H\textasciicircum{}0 of the internal hom into a K-injective complex computes Hom in the derived category}
\author{The Stacks Project Authors}
\date{stacks-project:sites-cohomology.tex:local-snapshot-20261011; GNU FDL 1.2 or later}
\begin{document}
\maketitle
\noindent\textit{Editorial note.} All delivered mathematical text is the author's own, copied from the pinned Stacks Project chapter file sites-cohomology.tex; no statement or proof step was rewritten, reordered or removed.\par\medskip
\noindent\textit{Editorial note.} The counted proof is the author's complete argument: the proof is written out in full: the displayed chain for the global statement and the localisation step through the restriction lemma for K-injective complexes.\par\medskip
\noindent\textit{Editorial note.} The label lemma-RHom-into-K-injective names the counted statement, and the proof environment immediately following it in the source is delivered verbatim as its proof.\par\medskip
\noindent\textit{Editorial note.} The counted statement is carried by the author's own lemma environment in the chapter source.\par\medskip
\noindent\textit{Editorial note.} The statement and proof are copied from the pinned chapter file; the three cited inputs are kept as external pointers with their canonical Stacks tags.\par\medskip
\noindent\textit{Allowed standard prerequisites (editorial).}\begin{itemize}
\item The cohomology of the internal hom complex as homotopy classes of maps. Same chapter, the equation identifying the cohomology of the internal hom complex with the homotopy category Hom group. Supplies the first equality of the displayed chain.
\item Localisation of K-injectivity to an open. Same chapter, the lemma stating that the restriction of a K-injective complex to a localisation is again K-injective. Supplies the second equality, since a K-injective representative computes Hom in the derived category.
\end{itemize}\medskip
\ifdefined\hypertarget\hypertarget{sourcepointer-1}{}\fi
\noindent\textit{Source reference.} Original equation 0A5X: sources/stacks-sites-cohomology/pinned-original-source.tex, source lines 8601--8601. Prerequisite pointer invoked at proof line(s) 8746; it names a standard fact of the same chapter that the counted argument uses, and every step of the counted proof built on it is retained verbatim in the delivered body.\par\medskip
\ifdefined\hypertarget\hypertarget{sourcepointer-2}{}\fi
\noindent\textit{Source reference.} Original equation 0A8Z: sources/stacks-sites-cohomology/pinned-original-source.tex, source lines 8608--8608. Prerequisite pointer invoked at proof line(s) 8750; it names a standard fact of the same chapter that the counted argument uses, and every step of the counted proof built on it is retained verbatim in the delivered body.\par\medskip
\ifdefined\hypertarget\hypertarget{sourcepointer-3}{}\fi
\noindent\textit{Source reference.} Original reference 08BS: sources/stacks-sites-cohomology/pinned-original-source.tex, source lines 4314--4314. Prerequisite pointer invoked at proof line(s) 8748; it names a standard fact of the same chapter that the counted argument uses, and every step of the counted proof built on it is retained verbatim in the delivered body.\par\medskip
\begin{lemma}
\label{lemma-RHom-into-K-injective}
Let $(\mathcal{C}, \mathcal{O})$ be a ringed site. Let $L$ and $M$
be objects of $D(\mathcal{O})$. Let $\mathcal{I}^\bullet$
be a K-injective complex of $\mathcal{O}$-modules representing $M$. Let
$\mathcal{L}^\bullet$ be a complex of $\mathcal{O}$-modules
representing $L$.
Then
$$
H^0(\Gamma(U, \SheafHom^\bullet(\mathcal{L}^\bullet, \mathcal{I}^\bullet))) =
\Hom_{D(\mathcal{O}_U)}(L|_U, M|_U)
$$
for all $U \in \Ob(\mathcal{C})$. Similarly,
$H^0(\Gamma(\mathcal{C},
\SheafHom^\bullet(\mathcal{L}^\bullet, \mathcal{I}^\bullet))) =
\Hom_{D(\mathcal{O})}(L, M)$.
\end{lemma}

\begin{proof}
We have
\begin{align*}
H^0(\Gamma(U, \SheafHom^\bullet(\mathcal{L}^\bullet, \mathcal{I}^\bullet)))
& =
\Hom_{K(\mathcal{O}_U)}(\mathcal{L}^\bullet|_U, \mathcal{I}^\bullet|_U) \\
& =
\Hom_{D(\mathcal{O}_U)}(L|_U, M|_U)
\end{align*}
The first equality is (\ref{equation-cohomology-hom-complex}).
The second equality is true because $\mathcal{I}^\bullet|_U$
is K-injective by Lemma \ref{lemma-restrict-K-injective-to-open}.
The proof of the last equation is similar except that it uses
(\ref{equation-global-cohomology-hom-complex}).
\end{proof}

\end{document}
